%------------------------------------------------------------------------------%
\section{O conceito de Ângulo}

Um \emph{ângulo}\index{Ângulo}
é determinado por uma abertura entre duas retas de mesma
origem. Para medirmos ângulos precisamos estabelecer o valor do ângulo total
(uma volta completa no sentido anti-horário) e os demais ângulos são medidos de
maneira proporcional. Existem duas formas usuais de medirmos ângulos: definindo
o ângulo total como 360 \emph{graus}\index{Graus}
ou como $2\pi$ \emph{radianos}\index{Radianos}.
Vejamos algumas medições em cada um dos dois sistemas.

%------------------------------------------------------------------------------%
\begin{figure}
  \centering
  \includegraphics[width=0.3\textwidth,trim={0 0 10,5cm 1cm},clip]{circ1}
  \qquad
  \includegraphics[width=0.3\textwidth,trim={0 0 10,5cm 1cm},clip]{circ2}
  \caption{Círculo 1 e 2}
\end{figure}
%------------------------------------------------------------------------------%

Considere um ângulo $\theta$. Denotamos por $\theta_g$ a medida desse ângulo em
graus e $\theta_r$ em radianos. Esses dois valores estão relacionados pela
equação. Isso quer dizer que podemos obter $\theta_g$ se conhecemos
$\theta_r$ ou vice versa.
\[
  \frac{360}{2\pi} = \frac{\theta_g}{\theta_r}.
\]
A unidade de medida padrão no ensino superior é radianos e para transformar
para radianos é comum usar a relação equivalente $\ang{180}=\pi$ rad,
conforme ilustra o exemplo que segue.

%------------------------------------------------------------------------------%
\begin{example}
  Converta $120^{\circ}$ para radianos.

  \solution
  Para fazermos essa conversão usamos que \ang{180} é equivalente
  a $\pi rad$ e montamos a seguinte regra de três.
  \[
    \begin{array}{c<{\quad}c}
      \text{graus} & \text{radianos} \\ \hline
      \ang{180}    & \pi \text{ rad} \\
      \ang{120}    & x
    \end{array}
  \]
  Assim,
  \begin{align*}
    180 x & = 120\pi                       \\
    x     & = \frac{2 \pi}{3} \text{ rad}.
  \end{align*}
  Logo, \ang{120} é equivalente a \(\frac{2 \pi}{3}\) rad.
\end{example}
%------------------------------------------------------------------------------%

Pode ser que seja necessário voltar de radianos para graus.
O exemplo seguinte ilustra essa situação.

%------------------------------------------------------------------------------%
\begin{example}
 Converta $\frac{5\pi}{3}$ radianos para graus.

 \solution
 Para converter de radianos para graus, basta usarmos a relação
 $\ang{180} = \pi$ rad. Assim,
 \[
   \frac{5\pi}{3}\text{ rad}
   = 5\,\frac{\ang{180}}{3}
   = \ang{300}.
 \]
\end{example}
%------------------------------------------------------------------------------%

